\documentclass[a4paper,12pt]{article}
\usepackage[T2A]{fontenc}
\usepackage[utf8]{inputenc}
\usepackage[russian]{babel}
\usepackage{amsmath,amssymb,mathtools}
\usepackage{helvet}
\renewcommand{\familydefault}{\sfdefault}
\usepackage{icomma}
\usepackage[left=19mm,right=19mm,top=20mm,bottom=20mm]{geometry}
\usepackage{microtype}
\usepackage{tikz}
\usetikzlibrary{calc,arrows.meta,positioning,shapes.geometric}
\usepackage{tabularx,booktabs,longtable}
\usepackage{enumitem}
\usepackage{fancyhdr}
\usepackage{parskip}
\usepackage{setspace}
\usepackage{xcolor}
\usepackage[hidelinks]{hyperref}
\definecolor{accent}{HTML}{2A6F6B}
\definecolor{darkaccent}{HTML}{174B48}
\definecolor{textgray}{HTML}{2F3437}
\definecolor{softgray}{HTML}{C8D8D6}
\color{textgray}
\onehalfspacing
\raggedbottom
\setlength{\parindent}{0pt}
\setlength{\headheight}{25pt}
\setlist[itemize]{leftmargin=7mm,topsep=3pt,itemsep=3pt}
\setlist[enumerate]{leftmargin=8mm,topsep=6pt,itemsep=9pt}
\pagestyle{fancy}
\fancyhf{}
\fancyhead[L]{\small Вклады и схемы погашения кредита}
\fancyhead[R]{\small 20.10.26}
\fancyfoot[C]{\small\thepage}
\renewcommand{\headrulewidth}{0.35pt}
\renewcommand{\headrule}{\hbox to\headwidth{\color{softgray}\leaders\hrule height \headrulewidth\hfill}}
\newcommand{\key}[1]{\textcolor{darkaccent}{\textbf{#1}}}
\newcommand{\rub}{\,\text{руб.}}
\newcommand{\mln}{\,\text{млн руб.}}
\newcommand{\fboxeq}[1]{\[\boxed{#1}\]}
\newcommand{\pct}{\,\%}
\begin{document}
\thispagestyle{empty}
\begin{center}
\vspace*{26mm}
{\Large\textcolor{accent}{\textbf{Чек-лист}}}\\[11mm]
{\huge\textbf{Финансовые задачи}}\\[8mm]
{\LARGE\textbf{Вклады и аннуитетный кредит}}\\[5mm]
{\large Сравнение способов погашения}\\[25mm]
{\large 20.10.26}\\[25mm]
{\large Лёвин Артём Александрович}\\[4mm]
{\large\textcolor{darkaccent}{эксклюзивно для Екатерины Гнедковой}}
\end{center}
\vfill
\begin{center}
\begin{tikzpicture}[x=1cm,y=1cm]
\draw[accent!65,line width=1.05pt]
(-7,0) .. controls (-4.5,0.8) and (-1.8,-0.7) .. (1.1,0.10)
.. controls (3.1,0.7) and (5.2,-0.55) .. (7,0.22);
\end{tikzpicture}
\end{center}
\clearpage

\section*{1. Как различать финансовые схемы}

Обозначим $S$ --- первоначальную сумму, $n$ --- количество лет, $r$ --- годовую ставку в процентах, $p=r/100$ --- ставку в долях единицы, $k=1+p$ --- коэффициент начисления процентов. В рассматриваемых задачах $n\in\mathbb N$, $n\ge1$, $S>0$, $p\ge0$. Проценты начисляются \key{раз в год} на имеющийся остаток перед очередным взносом или платёжом.

\begin{center}
\renewcommand{\arraystretch}{1.28}
\begin{tabularx}{\linewidth}{@{}p{37mm}X@{}}
\toprule
\textbf{Фраза из условия} & \textbf{Что означает}\\
\midrule
«Вклад ежегодно пополняли» & После начисления процентов на счёт вносится дополнительная сумма.\\
«Погасили кредит равными платежами» & \key{Аннуитет}: каждый период платят одинаковую сумму $x$.\\
«Долг уменьшался на одну и ту же сумму» & \key{Дифференцированное погашение}: одинаковые части основного долга; полный платёж уменьшается.\\
\bottomrule
\end{tabularx}
\end{center}

\key{Не путайте:} при аннуитете постоянен \emph{весь платёж}, при дифференцированной схеме --- \emph{часть платежа, возвращающая основной долг}.

\section*{2. Вклад с одинаковыми пополнениями}

Если вклад $S$ лежит $n$ лет под постоянную ставку, а после \emph{каждого из первых $n-1$ лет} (после начисления процентов) добавляют сумму $a$, то конечная сумма $F$ равна
\fboxeq{F=Sk^n+a\bigl(k^{n-1}+k^{n-2}+\cdots+k\bigr).}
Для снятия одной и той же суммы $a$ в эти моменты перед скобкой ставится \emph{минус}; сниматься может лишь сумма, имеющаяся на счёте.

\textbf{Почему именно такие степени?} Первый взнос начисляет проценты ещё $n-1$ раз, последний --- один раз. В финале $n$-го года добавочного взноса \emph{нет}. Если он предусмотрен условием, добавьте ещё одно $+a$.

\textbf{Пример.} Вклад $1\mln$ положили на $3$ года под $20\pct$ годовых. После первого и второго года вносили по $a\mln$. К концу третьего года стало $3{,}312\mln$. Тогда
\[
3{,}312=1\cdot1{,}2^3+a(1{,}2^2+1{,}2)
=1{,}728+2{,}64a,
\quad a=0{,}6.
\]
Значит, ежегодное пополнение составляло $600\,000\rub$.

\key{Если процент меняется по годам,} надёжнее составить таблицу: начальная сумма $\to$ начисленные проценты $\to$ пополнение (либо снятие) $\to$ новый остаток.

\clearpage
\section*{3. Аннуитетная схема: от таблицы к формуле}

Аннуитетный кредит погашается \key{равными платежами $x$} в конце каждого года, сразу после начисления процентов. Обозначим $B_j$ --- долг после $j$-го платежа. Тогда
\[
B_0=S,\qquad B_j=kB_{j-1}-x,\qquad B_n=0.
\]
Последнее равенство означает, что кредит погашен полностью.

Для трёх лет последовательное вычисление остатков даёт:
\[
\begin{aligned}
B_1&=Sk-x,\\
B_2&=(Sk-x)k-x=Sk^2-xk-x,\\
B_3&=(Sk^2-xk-x)k-x\\
   &=Sk^3-x(k^2+k+1)=0.
\end{aligned}
\]
Отсюда для трёх равных платежей:
\fboxeq{x=\frac{Sk^3}{k^2+k+1}.}
Аналогично для $n$ лет:
\fboxeq{Sk^n=x\bigl(k^{n-1}+k^{n-2}+\cdots+k+1\bigr).}
\fboxeq{x=\frac{Sk^n}{k^{n-1}+k^{n-2}+\cdots+k+1}.}

\textbf{Смысл формулы:} кредитная сумма $S$ накапливает проценты $n$ раз; каждый внесённый платёж уменьшает задолженность и тем самым влияет на дальнейшие начисления процентов.

\section*{4. Общая сумма, переплата и процент переплаты}
\enlargethispage{2\baselineskip}

После нахождения ежегодного платежа $x$ последовательно вычислите:
\[
\boxed{T=nx},\qquad
\boxed{Q=T-S=nx-S},\qquad
\boxed{q=\frac{Q}{S}\cdot100\pct}.
\]
Здесь $T$ --- сумма всех платежей; $Q$ --- переплата в рублях; $q$ --- \emph{процент} переплаты от суммы кредита. Годовая ставка $r$ и процент переплаты $q$ обычно \key{не совпадают}.

\textbf{Пример.} $S=210\,000\rub$, $n=2$, $r=10\pct$. Тогда $k=1{,}1$ и
\[
x=\frac{210\,000\cdot1{,}1^2}{1{,}1+1}=121\,000\rub.
\]
Проверка остатков: $B_1=231\,000-121\,000=110\,000$;
$B_2=121\,000-121\,000=0$. Поэтому
\[
T=242\,000\rub,\quad Q=32\,000\rub,\quad q\approx15{,}24\pct.
\]

\clearpage
\section*{5. Дифференцированная схема и сравнение}

При дифференцированном погашении ежегодно возвращается одинаковая часть \emph{основного долга} $S/n$. Проценты начисляются на остаток перед очередным платежом, поэтому полный платёж обычно убывает.

Для года с номером $j$ ($1\le j\le n$):
\[
P_j=\frac Sn+p\,\frac{S(n-j+1)}{n}.
\]
При постоянной ставке суммарные выплаты и переплата равны
\fboxeq{T_{\text{д}}=S+Sp\frac{n+1}{2},\qquad
Q_{\text{д}}=Sp\frac{n+1}{2}.}

\textbf{Пример сравнения из условия урока.} Кредит $100\,000\rub$ на $3$ года под $20\pct$ годовых.
\begin{center}
\renewcommand{\arraystretch}{1.25}
\begin{tabularx}{\linewidth}{@{}p{39mm}XX@{}}
\toprule
& \textbf{Дифференцированная} & \textbf{Аннуитетная} \\
\midrule
Ежегодный платёж & Убывает & Постоянный $x$\\
Общая выплата & $140\,000\rub$ & $\approx142\,417{,}58\rub$\\
Переплата & $40\,000\rub$ & $\approx42\,417{,}58\rub$\\
\bottomrule
\end{tabularx}
\end{center}
Аннуитетный платёж здесь равен $100\,000\cdot1{,}2^3/(1{,}2^2+1{,}2+1)\approx47\,472{,}53\rub$. Дифференцированная схема экономит $\approx2\,417{,}58\rub$ за весь срок.

\key{Важно:} при сравнении необходимо зафиксировать одну и ту же сумму кредита, ставку, срок и момент платежей. В реальном договоре могут быть комиссии и другие условия; в школьных задачах ориентируйтесь на условия, которые даны явно.

\section*{6. Типичные ошибки}
\begin{itemize}
\item Перепутать $r$ процентов и $p=r/100$: при $20\pct$ имеем $p=0{,}2$, $k=1{,}2$.
\item Применить формулу аннуитета к одинаковым частям \emph{основного долга}.
\item Не учесть момент пополнения вклада и лишний раз умножить взнос на $k$.
\item Считать проценты на исходный долг $S$ каждый год, игнорируя остаток.
\item Принять суммарные платежи $nx$ за переплату: обязательно вычтите $S$.
\item Округлять годовой платёж слишком рано: при сравнении используйте точное выражение и округляйте в конце.
\end{itemize}

\clearpage
\section*{Алгоритм: как решить задачу с равными платежами}
\vspace{3mm}
\begin{center}
\begin{tikzpicture}[
  >={Latex[length=2.5mm]},
  line width=.85pt,
  start/.style={ellipse,draw=darkaccent,fill=softgray!24,minimum width=47mm,minimum height=10mm,align=center,font=\normalsize},
  action/.style={rectangle,rounded corners=1.5mm,draw=darkaccent,minimum height=12mm,text width=65mm,inner sep=3mm,align=center,font=\small},
  short/.style={rectangle,rounded corners=1.5mm,draw=darkaccent,minimum height=14mm,text width=53mm,inner sep=3mm,align=center,font=\small},
  question/.style={diamond,aspect=2.7,draw=darkaccent,fill=softgray!18,inner xsep=3mm,inner ysep=2mm,text width=38mm,align=center,font=\small},
  link/.style={->,draw=darkaccent}
]
\node[start] (begin) at (0,0) {Дано: срок $n$, ставка $p$};
\node[action] (k) at (0,-2.1) {Найдите коэффициент $k=1+p$};
\node[action] (balance) at (0,-4.35) {Используйте $B_j=kB_{j-1}-x$,\\ $B_0=S$, $B_n=0$};
\node[question] (choose) at (0,-7.3) {Что требуется найти?};
\node[short] (pay) at (-3.8,-10.25) {Платёж $x$:\\[2pt]$\displaystyle x=\frac{Sk^n}{1+k+\cdots+k^{n-1}}$};
\node[short] (sum) at (3.8,-10.25) {Сумму кредита $S$:\\[2pt]$\displaystyle S=\frac{x(1+k+\cdots+k^{n-1})}{k^n}$};
\node[action] (total) at (0,-13.35) {Если требуется, найдите $T=nx$,\\ $Q=T-S$, $q=100Q/S\pct$};
\node[start] (end) at (0,-15.8) {Проверьте $B_n=0$};
\draw[link] (begin) -- (k);
\draw[link] (k) -- (balance);
\draw[link] (balance) -- (choose);
\draw[link] (choose.south west) -| node[pos=.22,above left,font=\small]{$x$} (pay.north);
\draw[link] (choose.south east) -| node[pos=.22,above right,font=\small]{$S$} (sum.north);
\draw[link] (pay.south) |- (total.west);
\draw[link] (sum.south) |- (total.east);
\draw[link] (total) -- (end);
\end{tikzpicture}
\end{center}
\clearpage

\section*{7. Тренировочный блок --- Путь А}
\textbf{10 заданий.} Решайте по порядку. В задачах о кредите проценты начисляются раз в год перед платежом, комиссии отсутствуют. Если не оговорено иное, платежи выполняются в конце года. Округляйте денежный ответ до копеек только в конце вычислений.

\begin{enumerate}
\item Определите тип каждой ситуации: (а) долг каждый год уменьшается на $30\,000\rub$; (б) все три платежа одинаковы; (в) вклад пополняют после каждого из первых двух лет. Назовите три схемы.

\item На вклад внесли $500\,000\rub$ под $10\pct$ годовых на два года. После первого года (после начисления процентов) добавили $200\,000\rub$. Сколько денег окажется на счёте в конце второго года?

\item На вклад внесли $1\mln$ на три года под $20\pct$ годовых. После первого и второго годов (после начисления процентов) вносили одинаковую сумму. Через три года на счёте стало $3{,}048\mln$. Найдите ежегодное пополнение.

\item Вклад $1\mln$ размещён на три года под $10\pct$ годовых. После первого и второго годов сразу после начисления процентов снимали по $100\,000\rub$. Найдите сумму на счёте в конце третьего года.

\item Кредит $150\,000\rub$ взят на три года под $12\pct$ годовых с равномерным погашением основного долга. Найдите все три платежа, общую сумму выплат и переплату.

\item При дифференцированном погашении кредита за пять лет общая переплата составила $15\pct$ от первоначальной суммы. Найдите постоянную годовую процентную ставку.

\item Кредит $210\,000\rub$ взят на два года под $10\pct$ годовых. По договору ежегодные платежи одинаковы. Найдите каждый платёж, общую сумму выплат и переплату.

\item Кредит $91\,000\rub$ взят на три года под $20\pct$ годовых и погашается равными годовыми платежами. Найдите размер одного платежа, переплату в рублях и процент переплаты (процент округлите до сотых).

\item Кредит на три года под $10\pct$ годовых погашается тремя равными платежами по $133\,100\rub$. Какова первоначальная сумма кредита?

\item Человек выбирает способ погашения кредита $100\,000\rub$ на три года под $20\pct$ годовых: равными платежами или равными частями основного долга. Какой способ даёт меньшую общую выплату и на сколько рублей?
\end{enumerate}
\clearpage

\section*{8. Ответы и самопроверка}
\small
\renewcommand{\arraystretch}{1.38}
\begin{longtable}{@{}p{11mm}p{147mm}@{}}
\toprule
\textbf{№} & \textbf{Ответ}\\
\midrule
\endhead
1 & (а) дифференцированная; (б) аннуитетная; (в) вклад с пополнениями.\\
2 & $825\,000\rub$ ($500\,000\cdot1{,}1^2+200\,000\cdot1{,}1$).\\
3 & $500\,000\rub$; $3{,}048=1{,}728+a(1{,}44+1{,}2)$, здесь $a$ в млн руб.\\
4 & $1\,100\,000\rub$; $1\,000\,000\cdot1{,}1^3-100\,000(1{,}1^2+1{,}1)$.\\
5 & $68\,000$, $62\,000$, $56\,000\rub$; всего $186\,000\rub$, переплата $36\,000\rub$.\\
6 & $5\pct$ годовых: $0{,}15=p(5+1)/2$.\\
7 & $x=121\,000\rub$; $T=242\,000\rub$; $Q=32\,000\rub$.\\
8 & $x=43\,200\rub$; $Q=38\,600\rub$; $q\approx42{,}42\pct$.\\
9 & $S=331\,000\rub$: $S=133\,100(1+1{,}1+1{,}1^2)/1{,}1^3$.\\
10 & Выгоднее \textbf{дифференцированная схема}: $T_{\text{д}}=140\,000\rub$, $T_{\text{а}}\approx142\,417{,}58\rub$; экономия $\approx2\,417{,}58\rub$.\\
\bottomrule
\end{longtable}
\normalsize
\vspace{3mm}
\textbf{Контрольные вопросы}
\begin{itemize}
\item В какой схеме постоянен сам платёж, а в какой --- часть основного долга?
\item Почему равный ежегодный платёж при аннуитете приводит к остатку $B_n=0$?
\item Как по найденному платежу вычислить общую выплату, переплату и процент переплаты?
\item Почему при расчёте вклада нужно уточнять \emph{момент} каждого пополнения?
\end{itemize}
\end{document}
